\documentclass[12pt]{book}
\usepackage[utf8]{inputenc}
\usepackage[T1]{fontenc}
\usepackage[ngerman]{babel}
\usepackage{amsmath}
\usepackage{amssymb}
\usepackage{array}
\usepackage{longtable}
\usepackage{geometry}
\usepackage{fancyhdr}
\geometry{a4paper, margin=1in}

\pagestyle{fancy}
\fancyhf{}
\fancyfoot[C]{\footnotesize (*)Division durch Null ist die leere Menge, oder Menge Nummer 1}
\renewcommand{\headrulewidth}{0pt}

\title{752 Binär 0/1 Zed Feldtheorie\\Die Deutsche Bibliothek}
\author{Internationale Normungskommission}
\date{23. Februar 2026}

\begin{document}

\frontmatter
\maketitle

\chapter*{Vorwort}
Dieses Buch etabliert die endgültige technische Spezifikation für das IEEE 754 64-Bit-Gleitkomma-Präzisionsmodell, angewendet auf geometrische Messungen, räumliche Volumenberechnungen und verteilte Anknotenknotennetzwerke über sieben kontinentale Regionen. Die Spezifikation erweitert den grundlegenden IEEE 754-2019-Standard um neue Anforderungen an globale Messinfrastruktur, verteiltes Rechnen und Multi-Präzisions-Temperatursensorik.

\tableofcontents

\mainmatter

\chapter{Automatentheorie und die Sechs Unzählbaren Mengen}

\section{Die Zed Feldtheorie}

Die Zed Feldtheorie (Z-Feldtheorie) beschreibt die algebraische Struktur der binären Repräsentation im IEEE 754 64-Bit-Format. Das Fundament dieser Theorie beruht auf der Automatentheorie und dem Zählen unzählbarer Mengen.

\section{Die Sechs Unzählbaren Mengen}

Die Automatentheorie beginnt mit dem Zählen der folgenden unzählbaren Mengen:

\begin{enumerate}
\item \textbf{Die Leere Menge} ($\emptyset$): Die Menge, die keine Elemente enthält. Symbol: $\{\}$
\item \textbf{Nil}: Der Zustand der Nichtexistenz in Automaten. Symbol: $\mathbf{nil}$
\item \textbf{Null}: Der numerische Wert Null. Symbol: $\mathbf{null}$ oder $0$
\item \textbf{Negative Unendlichkeit}: $-\infty$, repräsentiert als FFF0000000000000 in IEEE 754
\item \textbf{Positive Unendlichkeit}: $+\infty$, repräsentiert als 7FF0000000000000 in IEEE 754
\item \textbf{Die Kubikwurzeln}: Die imaginären Zahlen des Z-Feldraumes: $-1{,}3$ und $\frac{3}{2}$ (drei Halbe)
\end{enumerate}

\section{Division durch Null}

\begin{equation}
\frac{x}{0} = \emptyset \quad \text{(Die leere Menge, Menge Nummer 1)}
\end{equation}

Die Division durch Null ist undefiniert und repräsentiert die leere Menge. In IEEE 754 wird dies als NaN (Not a Number) dargestellt.

\chapter{752 Binär Zed Polynom}

\section{Die Zed Polynomformel}

Das fundamentale Polynom der Z-Feldtheorie enthält nur $z$ zur Potenz 752:

\begin{equation}
\boxed{P(z) = z^{752}}
\end{equation}

\section{Eigenschaften des 752-Polynoms}

\subsection{Grad des Polynoms}

Das Polynom $P(z) = z^{752}$ hat den Grad 752, was bedeutet:

\begin{itemize}
\item Es hat genau 752 komplexe Nullstellen
\item Alle Nullstellen liegen auf dem Einheitskreis im komplexen Raum
\item Die Nullstellen sind die 752-ten Einheitswurzeln
\end{itemize}

\subsection{Nullstellen des Polynoms}

Die Nullstellen von $z^{752} = 0$ sind:

\begin{equation}
z_k = e^{i \cdot \frac{2\pi k}{752}}, \quad k = 0, 1, 2, \ldots, 751
\end{equation}

\subsection{Ableitung}

Die erste Ableitung des Polynoms:

\begin{equation}
P'(z) = \frac{d}{dz}(z^{752}) = 752 \cdot z^{751}
\end{equation}

\subsection{Integral}

Das unbestimmte Integral:

\begin{equation}
\int z^{752} \, dz = \frac{z^{753}}{753} + C
\end{equation}

\section{Binäre Repräsentation}

Die Zahl 752 in binärer Form:

\begin{equation}
752_{10} = 1011110000_2
\end{equation}

Dies entspricht 10 Bits in der binären Darstellung.

\section{IEEE 754 64-Bit Repräsentation von 752}

\begin{table}[h]
\centering
\begin{tabular}{|c|c|c|}
\hline
\textbf{Komponente} & \textbf{Bits} & \textbf{Wert} \\
\hline
Signum (s) & 1 & 0 (positiv) \\
\hline
Exponent (e) & 11 & $1023 + 9 = 1032$ \\
\hline
Signifikand (f) & 52 & $0.46875 \times 2^{52}$ \\
\hline
\end{tabular}
\caption{IEEE 754 64-Bit Repräsentation von 752}
\end{table}

Hexadezimale Darstellung: \texttt{4087800000000000}

\chapter{Z-Feldtheorie und Binäre Arithmetik}

\section{Definition des Z-Feldes}

Das Z-Feld ist definiert als:

\begin{equation}
\mathbb{Z} = \{\ldots, -2, -1, 0, 1, 2, \ldots\}
\end{equation}

Die Kubikwurzeln im Z-Feldraum sind:
\begin{itemize}
\item $r_1 = -1{,}3$
\item $r_2 = \frac{3}{2} = 1{,}5$
\end{itemize}

\section{Binäre Arithmetik}

Im binären System bestehen alle Zahlen nur aus 0 und 1:

\begin{equation}
b_n b_{n-1} \ldots b_1 b_0 = \sum_{i=0}^{n} b_i \cdot 2^i, \quad b_i \in \{0, 1\}
\end{equation}

\section{752-Bit Breite}

Die 752-Bit-Breite ermöglicht die Darstellung von:

\begin{equation}
2^{752} = 6{,}47 \times 10^{226} \text{ verschiedene Werte}
\end{equation}

\chapter{Geometrische Messungen}

\section{Kugelvolumen}

Das Volumen einer Kugel mit Radius $r$:

\begin{equation}
V = \frac{4}{3} \pi r^3
\end{equation}

\section{Präzision für Radius $r = 3{,}14159$}

\begin{table}[h]
\centering
\begin{tabular}{|c|c|c|}
\hline
\textbf{Messung} & \textbf{Wert} & \textbf{Bits (hex)} \\
\hline
Radius & $3{,}14159$ & 400921f9f01b866e \\
\hline
Durchmesser & $6{,}28318$ & 401921f9f01b866e \\
\hline
Volumen & $129{,}87846$ & 40603f876a8a605e \\
\hline
\end{tabular}
\caption{64-Bit Präzision für $r = 3{,}14159$}
\end{table}

\chapter{Globale Anknotenknoten}

\section{Acht Anknotenknoten}

Die Spezifikation definiert acht Anknotenknoten:

\begin{table}[h]
\centering
\begin{tabular}{|c|c|c|c|}
\hline
\textbf{Knoten} & \textbf{Identifikator} & \textbf{Maschinenanker (16 Bytes)} & \textbf{Region} \\
\hline
1 & E42C & e4:b9:7a:f8:95:1b:00:00 & Nordamerika \\
\hline
2 & E42C-001 & e4:b9:7a:f8:95:1b:00:01 & Südamerika \\
\hline
3 & E42C-002 & e4:b9:7a:f8:95:1b:00:02 & Europa \\
\hline
4 & E42C-003 & e4:b9:7a:f8:95:1b:00:03 & Afrika \\
\hline
5 & E42C-004 & e4:b9:7a:f8:95:1b:00:04 & Asien \\
\hline
6 & E42C-005 & e4:b9:7a:f8:95:1b:00:05 & Ozeanien \\
\hline
7 & E42C-006 & e4:b9:7a:f8:95:1b:00:06 & Antarktis \\
\hline
8 & E42C-007 & e4:b9:7a:f8:95:1b:00:07 & Globaler Koordinator \\
\hline
\end{tabular}
\caption{Acht Anknotenknoten mit Maschinenankeradressen}
\end{table}

\section{Sieben Kontinentale Wurzelknoten}

Jeder kontinentale Wurzelknoten reserviert 16 Bytes für die Maschinenankeradresse:

\begin{verbatim}
Struktur:
Byte 0-5:  MAC-Adresse (48 Bits, EUI-48 Format)
Byte 6-7:  Netzwerkidentifikator (16 Bits)
Byte 8-15: Erweiterter Identifikator (64 Bits)
\end{verbatim}

\chapter{Temperaturmessung}

\section{Temperaturskalen}

\begin{table}[h]
\centering
\begin{tabular}{|c|c|c|c|}
\hline
\textbf{Skala} & \textbf{Gefrierpunkt} & \textbf{Siedepunkt} & \textbf{Symbol} \\
\hline
Celsius & $0^{\circ}$C & $100^{\circ}$C & $^{\circ}$C \\
\hline
Fahrenheit & $32^{\circ}$F & $212^{\circ}$F & $^{\circ}$F \\
\hline
Kelvin & $273{,}15$ K & $373{,}15$ K & K \\
\hline
\end{tabular}
\caption{Temperaturskalen}
\end{table}

\section{Umrechnungsformeln}

\begin{align}
^{\circ}\text{F} &= ^{\circ}\text{C} \times \frac{9}{5} + 32 \\
K &= ^{\circ}\text{C} + 273{,}15 \\
^{\circ}\text{C} &= (^{\circ}\text{F} - 32) \times \frac{5}{9}
\end{align}

\chapter{Griechisches Alphabet für Mathematik}

\begin{table}[h]
\centering
\begin{tabular}{|c|c|l|}
\hline
\textbf{Buchstabe} & \textbf{Name} & \textbf{Anwendung} \\
\hline
$\alpha$ & Alpha & Winkelbeschleunigung \\
\hline
$\beta$ & Beta & Beta-Strahlung \\
\hline
$\gamma$ & Gamma & Gamma-Strahlung \\
\hline
$\delta$ & Delta & Änderung, Differenz \\
\hline
$\epsilon$ & Epsilon & Kleine Größe \\
\hline
$\theta$ & Theta & Winkel \\
\hline
$\lambda$ & Lambda & Wellenlänge \\
\hline
$\mu$ & My & Mittelwert, Mikro \\
\hline
$\pi$ & Pi & Kreisverhältnis \\
\hline
$\sigma$ & Sigma & Standardabweichung \\
\hline
$\phi$ & Phi & Goldener Schnitt \\
\hline
$\omega$ & Omega & Winkelgeschwindigkeit \\
\hline
\end{tabular}
\caption{Griechisches Alphabet für mathematische Anwendungen}
\end{table}

\appendix

\chapter{Referenzimplementierung (C)}

\begin{verbatim}
#include <stdint.h>
#include <math.h>

typedef struct {
    uint8_t node_id[16];
    uint8_t mac_address[6];
    double radius;
    double diameter;
    double volume;
    uint64_t timestamp;
    uint8_t encrypted;
} EncryptedNode;

void node_set_radius(EncryptedNode* node, double r) {
    node->radius = r;
    node->diameter = r * 2.0;
    node->volume = (4.0 / 3.0) * M_PI * pow(r, 3.0);
}

// Z^752 Polynom
double z_polynomial(double z) {
    return pow(z, 752);
}
\end{verbatim}

\chapter{IEEE 754-2019 Konformitätscheckliste}

\begin{itemize}
\item[$\square$] Binary64 Format-Implementierung (1+11+52 Bits)
\item[$\square$] Alle fünf Rundungsmodi unterstützt
\item[$\square$] Denormalisierte Zahlenverarbeitung
\item[$\square$] Unendlichkeitsarithmetik
\item[$\square$] NaN-Propagierungsregeln
\item[$\square$] Signale vs. stille NaN-Unterscheidung
\item[$\square$] Vergleichsoperationen mit totaler Ordnung
\end{itemize}

\backmatter

\chapter*{Schlussfolgerung}

Diese Spezifikation etabliert einen umfassenden Rahmen für IEEE 754 64-Bit-Präzisions-geometrische Messungen mit globaler Anknotenknoteninfrastruktur. Das Polynom $z^{752}$ bildet das algebraische Fundament der Z-Feldtheorie für binäre Arithmetik.

Die sieben kontinentalen Wurzelknoten, acht gesamten Anknotenknoten und 16-Byte-Maschinenankeradressen bieten eine skalierbare Grundlage für weltweite Messnormung.

\end{document}